\chapter{Results and Evaluation}\label{ch:results}

This chapter presents the measurement results. Concrete diagrams and tables are generated
from the measurement runs and included from the (language-neutral) \texttt{tikz/} and
\texttt{tabellen/} directories; see the placeholders below. At the time of writing, the
methodology and the evaluation pipeline are in place and verified on sample data; the
real-hardware runs are pending (Section~\ref{sec:limitations}).

\section{Evaluation Pipeline}\label{sec:eval-pipeline}

Per series run, the \texttt{messung\_driver} writes
\texttt{Code/\_runs/<date>/<spec\_id>/} with \texttt{measurements.csv} (one row per
permutation: cycles, L1/L2/L3 misses, branches, throughput), \texttt{measurements.json},
and optional binary records. The evaluation uses \texttt{binary\_to\_csv} (deserialize and
enrich with profile id, paper ref, axes), \texttt{csv\_to\_latex} (booktabs tables with
per-profile fact sheets), and \texttt{diagram\_generator} (A4-aware TikZ bar charts,
scatter plots, heatmaps). The manuscript consumes these via
\texttt{\textbackslash input\{tikz/<spec\_id>/\ldots\}} and
\texttt{\textbackslash input\{tabellen/<spec\_id>\ldots\}}.

\section{Series A: PRT-ART vs.\ State of the Art}\label{sec:result-a}

% After the first messung_driver run, insert e.g.:
%   \input{tikz/A_full/cycles_per_permutation.tikz}
%   \input{tabellen/A_full_summary.tex}
The hypotheses H1--H3 (Section~\ref{sec:hypotheses}) are evaluated here once the
measurement data is available.

\section{Series B: Cache-Engine Permutations}\label{sec:result-b}

% \input{tikz/B_cache_engine_perms/heatmap.tikz}

\section{Series C: Full Join}\label{sec:result-c}

% \input{tikz/C_merge/comparison.tikz}

\section{Axis Sensitivity Analysis}\label{sec:sensitivity}

Using the per-permutation records, this section quantifies which axis contributes most to
the performance variance --- the basis for recommending default configurations per
workload class.

\section{Discussion}\label{sec:discussion}

The results are discussed against the leitmotif of the thesis: the transition from
heuristic to informed cache management, i.e.\ whether telemetry-driven, automatically
adapted configurations outperform statically chosen ones.
